\section{Conclusiones}
\label{sec:conclusions}

\begin{itemize}
  \item El dengue en el Perú crece por brotes cada vez más grandes y se concentra en la costa norte y la selva.
        Los distritos en emergencia por lluvias y con menos horas de agua al día tuvieron más incidencia, y en la
        costa los casos aumentan unos dos meses después de la lluvia. No se encontró relación entre la gravedad
        de los casos y la cobertura de establecimientos de salud.
  \item Sobre la pregunta central: Spark convierte las operaciones en un logical plan; Catalyst lo optimiza
        (pruning, pushdown, joins reescritos); el physical plan decide cómo ejecutar cada operador; y una action
        lo divide en stages en cada shuffle y ejecuta cada stage como tasks en paralelo.
  \item La API no cambia la ejecución: la misma consulta en DataFrame API y en Spark SQL llega al mismo plan.
  \item El diseño de los datos se refleja en el plan: el particionado por año y el formato Parquet reducen lo que
        se lee, y el broadcast de las tablas pequeñas evita shuffles.
  \item El tipo de join define el resultado, no solo el costo: con filas repetidas, un inner join multiplicó los
        casos, mientras que un left semi join los contó una sola vez.
\end{itemize}
